% Complete reviewed Unit 1; editorial master: chapters/ch01/answer-key.md
% Markdown SHA-256: 3946fa37fbdfff0ee9a3e5a153bd5866d26bcf7c51d5dfb29a52b046e7267d97
% This TeX is hand-maintained after the initial one-use assembly.
\newcommand{\runningtitle}{Unit 1. Answer key}
% Layout inherited from the reviewed previous course; set a running title per unit.
% This is a small hand-maintained TeX layout, not a conversion pipeline.
\RequirePackage{fix-cm}
\documentclass[11pt,twoside]{article}
\usepackage{fontspec}
\setmainfont{Cambria}[
  Renderer=HarfBuzz,
  BoldFont=cambriab.ttf,
  ItalicFont=cambriai.ttf,
  BoldItalicFont=cambriaz.ttf
]
\setmonofont{lmmono10-regular.otf}[BoldFont=lmmonolt10-bold.otf]
\usepackage[paperwidth=165mm,paperheight=241mm,inner=21mm,outer=19mm,
  top=22mm,bottom=20mm,headheight=14pt,headsep=7mm,footskip=10mm]{geometry}
\usepackage{xcolor,array,booktabs,longtable,tabularx}
\usepackage{enumitem,fancyhdr}
\usepackage[unicode,hidelinks]{hyperref}
\definecolor{ink}{HTML}{000000}
\definecolor{muted}{HTML}{4C5558}
\setlength{\parindent}{0pt}
\setlength{\parskip}{3.5pt}
\AtBeginDocument{\fontsize{11}{14.3}\selectfont}
\setlength{\tabcolsep}{4pt}
\renewcommand{\arraystretch}{1.10}
\setlist[itemize]{leftmargin=1.3em,label=\textbullet,itemsep=2pt,topsep=3pt,parsep=0pt}
\setlist[enumerate]{leftmargin=1.8em,itemsep=3pt,topsep=4pt,parsep=0pt}
\widowpenalty=10000
\clubpenalty=10000
\emergencystretch=1em
\raggedbottom
\pagestyle{fancy}
\fancyhf{}
\providecommand{\runningtitle}{Old English}
\fancyhead[LE,RO]{\small\itshape \runningtitle}
\fancyhead[LO,RE]{\small Old English}
\fancyfoot[LE,RO]{\small\thepage}
\renewcommand{\headrulewidth}{0.3pt}
\newcommand{\booktitle}[2]{
  {\fontsize{9}{11}\selectfont\color{muted}#1\par}
  \vspace{3pt}
  {\fontsize{20}{24}\selectfont\color{ink}#2\par}
  \vspace{6pt}
}
\newcommand{\source}[1]{\par{\footnotesize\color{muted}#1\par}}
\newcommand{\readingtext}[1]{
  \par\addvspace{5pt}
  \begingroup\leftskip=1.5em
    \fontsize{12}{16}\selectfont #1\par
  \endgroup\addvspace{5pt}
}
\makeatletter
\renewcommand{\section}{\@startsection{section}{1}{\z@}
  {-2.4ex \@plus -0.6ex \@minus -0.2ex}{0.8ex \@plus 0.2ex}
  {\normalfont\fontsize{12}{15}\selectfont\bfseries\color{ink}}}
\renewcommand{\subsection}{\@startsection{subsection}{2}{\z@}
  {-1.8ex \@plus -0.5ex \@minus -0.2ex}{0.5ex \@plus 0.2ex}
  {\normalfont\fontsize{11}{14}\selectfont\bfseries\color{ink}}}
\makeatother

\hyphenpenalty=10000
\exhyphenpenalty=10000
\emergencystretch=3em
\newcommand{\itempara}[2]{\par\addvspace{3pt}\begingroup\raggedright\noindent\hangindent=1.8em\hangafter=1\makebox[1.8em][l]{#1.}#2\par\endgroup}
\newcommand{\continuation}[1]{\par\begingroup\raggedright\leftskip=1.8em #1\par\endgroup}
\begin{document}
\raggedright
\booktitle{UNIT 1}{Answer key}

\section*{Reading the models}

The models use the forms selected in the lesson. A different grammatical answer may still convey the required meaning. Distinguish a mistake in Old English from failure to practise an explicitly requested form or word. In the supplied frames, the verb is given: assessment concerns the noun phrases and their cases, not recall of a verb paradigm.

For isolated Cook phrases, the head is nominative. Preserve singular and plural even where English treats the referent differently. Each article agrees with its own noun. A nominative/accusative form that happens to be identical is not evidence that the two syntactic roles are interchangeable.

% shared key note: KEY-NOMINAL-COOK
The source permits \textbf{ðām} beside \textbf{ðǣm}, \textbf{ðǣra} beside \textbf{ðāra}, and \textbf{sind/sint} beside \textbf{sindon}. Likewise, \textbf{fugolas/morgenas} and other appropriate unsyncopated forms, shorter \textbf{wætru}, bare singular \textbf{trēow}, strong feminine N/A plural \textbf{-e}, and singular oblique \textbf{-a} of the \textbf{-ung} group can be valid. When a task explicitly asks for the selected table, identify such an answer as an alternative rather than the requested model. Retain grammatical agreement throughout.

The selected \textbf{ǣfen} and \textbf{lyft} are neuter, and \textbf{sǣ} is masculine. Other evidenced genders are not intrinsically erroneous, but changing gender requires a consistent paradigm and agreement and does not meet a direction to use the selected form. Similarly, \textbf{nīeten} is a valid animal word, while the directions here choose \textbf{dēor} for “animal/beast” and \textbf{nīeten} for “cattle.” \textbf{Wæstm} has short \textbf{æ}; neither the erroneous printed macron in Cook nor \textbf{a} in its plural is the target. The starred \textbf{*sǣwa} is never required.

\section*{11. Forms and agreement}

\par\noindent\begin{minipage}{\linewidth}\raggedright
% item: CH01-F01
% key: KEY-CH01-F01
\itempara{1}{\textbf{Fisc with the article:}}

\par\begingroup
\setlength{\LTleft}{0pt}\setlength{\LTright}{\fill}
\noindent
\begin{tabular}{@{}>{\raggedright\arraybackslash}p{35mm}>{\raggedright\arraybackslash}p{40mm}>{\raggedright\arraybackslash}p{40mm}@{}}
\toprule
\textbf{Case} & \textbf{Singular} & \textbf{Plural} \\
\midrule
N & se fisc & ðā fiscas \\
G & ðæs fisces & ðāra fisca \\
D & ðǣm fisce & ðǣm fiscum \\
A & ðone fisc & ðā fiscas \\
\bottomrule
\end{tabular}
\par\endgroup

The target is the complete masculine paradigm with agreement. Singular nominative and accusative \textbf{fisc} coincide, but \textbf{se/ðone} distinguish them. Plural \textbf{-as} belongs to N/A, not to the genitive or dative.

\end{minipage}\par

\par\noindent\begin{minipage}{\linewidth}\raggedright
% item: CH01-F02
% key: KEY-CH01-F02
\itempara{2}{N/A plural \textbf{hofu; sǣd}. The short stem \textbf{hof-} takes \textbf{-u}; the long vowel of \textbf{sǣd} makes its stem heavy, so N/A plural adds no ending. Both cases have the same form within each neuter number. This is the stated ordinary pattern, not a rule to impose on \textbf{trēo} or the disyllabic nouns.}

\end{minipage}\par

\par\noindent\begin{minipage}{\linewidth}\raggedright
% item: CH01-F03
% key: KEY-CH01-F03
\itempara{3}{\textbf{Daga; hwala; męnn; nihta}, respectively. The first two use the plural \textbf{a} stem and genitive plural \textbf{-a}. \textbf{Męnn} is the exceptional dative singular of \textbf{mann}. \textbf{Nihta} is genitive plural even though the N/A plural \textbf{niht} has no ending.}

\end{minipage}\par

\par\noindent\begin{minipage}{\linewidth}\raggedright
% item: CH01-F04
% key: KEY-CH01-F04
\itempara{4}{\textbf{Trēowes, trēowe, trēowu; tācnes, tācnu; wæteres, wæteru; ǣfenu, nīetenu.} Every ending of \textbf{trēo} follows \textbf{w}. \textbf{Tācen} loses its middle \textbf{e}; the selected forms of the other three disyllables retain theirs. Dropping \textbf{-u} solely because the stem contains a long vowel misapplies the ordinary monosyllabic rule.}

\end{minipage}\par

\par\noindent\begin{minipage}{\linewidth}\raggedright
% item: CH01-F05
% key: KEY-CH01-F05
\itempara{5}{\textbf{Steorra} and the singular of \textbf{eorðe}:}

\par\begingroup
\setlength{\LTleft}{0pt}\setlength{\LTright}{\fill}
\noindent
\begin{tabular}{@{}>{\raggedright\arraybackslash}p{32mm}>{\raggedright\arraybackslash}p{28mm}>{\raggedright\arraybackslash}p{28mm}>{\raggedright\arraybackslash}p{28mm}@{}}
\toprule
\textbf{Case} & \textbf{steorra singular} & \textbf{steorra plural} & \textbf{eorðe singular} \\
\midrule
N & steorra & steorran & eorðe \\
G & steorran & steorrena & eorðan \\
D & steorran & steorrum & eorðan \\
A & steorran & steorran & eorðan \\
\bottomrule
\end{tabular}
\par\endgroup

The target is weak inflection: do not replace weak genitive singular \textbf{-an} with strong masculine \textbf{-es}. Plural genitive \textbf{-ena} and dative \textbf{-um} remain distinct from \textbf{-an}.

\end{minipage}\par

\par\noindent\begin{minipage}{\linewidth}\raggedright
% item: CH01-F06
% key: KEY-CH01-F06
\itempara{6}{\textbf{Wyrt, wyrte; tīd, tīde; gesceafte.} The order is accusative then genitive for the first two nouns. Their selected accusatives have no \textbf{-e}; the ordinary strong feminine accusative \textbf{gesceafte} does. A single feminine rule without these lexical exceptions is insufficient.}

\end{minipage}\par

\par\noindent\begin{minipage}{\linewidth}\raggedright
% item: CH01-F07
% key: KEY-CH01-F07
\itempara{7}{\textbf{Fæstnesse; andlīcnessa; gelīcnessa; gegaderungum.} Retain \textbf{e} and double \textbf{s} before the \textbf{-nes} nouns' endings. Genitive and nominative plural can share \textbf{-a}; identify the requested case from the task, not the ending alone.}

\end{minipage}\par

\par\noindent\begin{minipage}{\linewidth}\raggedright
% item: CH01-F08
% key: KEY-CH01-F08
\itempara{8}{\textbf{Ðȳ} and \textbf{ðon} are instrumental singular forms, both masculine and neuter. The feminine and plural have no separate instrumental forms in this table. Recognizing the cells is sufficient; no instrumental construction is being tested.}

\end{minipage}\par

\subsection*{The copula}

\par\noindent\begin{minipage}{\linewidth}\raggedright
% item: CH01-COP01
% key: KEY-CH01-COP01
\itempara{1}{First person singular: \textbf{eom / wæs} — present / preterite.}

\end{minipage}\par

\par\noindent\begin{minipage}{\linewidth}\raggedright
% item: CH01-COP02
% key: KEY-CH01-COP02
\itempara{2}{Second person singular: \textbf{eart / wǣre} — present / preterite.}

\end{minipage}\par

\par\noindent\begin{minipage}{\linewidth}\raggedright
% item: CH01-COP03
% key: KEY-CH01-COP03
\itempara{3}{Third person singular: \textbf{is / wæs} — present / preterite.}

\end{minipage}\par

\par\noindent\begin{minipage}{\linewidth}\raggedright
% item: CH01-COP04
% key: KEY-CH01-COP04
\itempara{4}{First person plural: \textbf{sindon / wǣron} — present / preterite.}

\end{minipage}\par

\par\noindent\begin{minipage}{\linewidth}\raggedright
% item: CH01-COP05
% key: KEY-CH01-COP05
\itempara{5}{Second person plural: \textbf{sindon / wǣron} — present / preterite.}

\end{minipage}\par

\par\noindent\begin{minipage}{\linewidth}\raggedright
% item: CH01-COP06
% key: KEY-CH01-COP06
\itempara{6}{Third person plural: \textbf{sindon / wǣron} — present / preterite.}

The forms are indicative. \textbf{Wǣre} is required for second-person singular preterite in this paradigm; \textbf{wæs} serves first and third singular. Each plural person has \textbf{sindon / wǣron}. Different persons may share a form without ceasing to be separate paradigm cells.

\end{minipage}\par

\section*{12. Case in sentences}

The following preserve the order of the printed frames. Other word orders can be grammatical, but changing the frame is not needed. An English “in” phrase has the location specified in the prompt; no motion construction is intended.

\par\noindent\begin{minipage}{\linewidth}\raggedright
% item: CH01-CASE01
% key: KEY-CH01-CASE01
\itempara{1}{\textbf{Se mann geseah ðone fugol.} — The man saw the bird. Subject se mann is nominative; object ðone fugol is accusative. Both noun spellings themselves could serve either case.}

\end{minipage}\par

\par\noindent\begin{minipage}{\linewidth}\raggedright
% item: CH01-CASE02
% key: KEY-CH01-CASE02
\itempara{2}{\textbf{Se fugol geseah ðone mann.} — The bird saw the man. The roles are reversed from item 1. Change the articles with the roles; do not merely reverse the English gloss.}

\end{minipage}\par

\par\noindent\begin{minipage}{\linewidth}\raggedright
% item: CH01-CASE03
% key: KEY-CH01-CASE03
\itempara{3}{\textbf{God gescēop ðā heofonas and ðā eorðan.} — God created the heavens and the earth. The coordinated objects are masculine accusative plural and feminine accusative singular. Their articles coincide as ðā. This task explicitly asks for masculine heofon, unlike the source reading’s supplied weak singular heofonan.}

\end{minipage}\par

\par\noindent\begin{minipage}{\linewidth}\raggedright
% item: CH01-CASE04
% key: KEY-CH01-CASE04
\itempara{4}{\textbf{Se mann geseah ðā steorran.} — The man saw the stars. Steorran is weak accusative plural here; ðā agrees with the plural object, not with the singular subject.}

\end{minipage}\par

\par\noindent\begin{minipage}{\linewidth}\raggedright
% item: CH01-CASE05
% key: KEY-CH01-CASE05
\itempara{5}{\textbf{Se mann geseah ðone fisc on ðǣm sǣ.} — The man saw the fish in the sea. The fish is accusative; sǣ is masculine dative singular after locative on. Sea singular nominative and dative happen to coincide.}

\end{minipage}\par

\par\noindent\begin{minipage}{\linewidth}\raggedright
% item: CH01-CASE06
% key: KEY-CH01-CASE06
\itempara{6}{\textbf{God gedǣlde ðæt lēoht fram ðǣm ðīestrum.} — God separated the light from the darkness. Light is the accusative object; darkness is dative plural after fram. This adapts the reading by using God in place of hē and omitting introductory and.}

\end{minipage}\par

\par\noindent\begin{minipage}{\linewidth}\raggedright
% item: CH01-CASE07
% key: KEY-CH01-CASE07
\itempara{7}{\textbf{Se fugol geseah ðæt sǣd on ðæs mannes hofe.} — The bird saw the seed in the man’s dwelling. Seed is accusative; hofe is dative after on. The internal dependent ðæs mannes remains genitive.}

\end{minipage}\par

\par\noindent\begin{minipage}{\linewidth}\raggedright
% item: CH01-CASE08
% key: KEY-CH01-CASE08
\itempara{8}{\textbf{Se fugol geseah ðone mann on ðǣm hofe.} — The bird saw the man in the dwelling. Mann is accusative; hofe is dative singular. Here the article belongs directly to hofe, unlike item 7.}

\end{minipage}\par

\par\noindent\begin{minipage}{\linewidth}\raggedright
% item: CH01-CASE09
% key: KEY-CH01-CASE09
\itempara{9}{\textbf{God sealde ðǣm męnn męte.} — God gave food to the man. The recipient ðǣm męnn is dative singular, with the changed-vowel mann form; the thing given, męte, is accusative singular. No preposition is required for the recipient.}

\end{minipage}\par

\par\noindent\begin{minipage}{\linewidth}\raggedright
% item: CH01-CASE10
% key: KEY-CH01-CASE10
\itempara{10}{\textbf{Se mann sealde ðǣm fugle sǣd.} — The man gave seed to the bird. Se mann is nominative, ðǣm fugle dative recipient, and sǣd accusative thing given. Bird dative singular is fugle, not fugol or fuglas.}

\end{minipage}\par

\clearpage
\section*{13. Cook's noun-phrase exercises}

\subsection*{Group I, Exercise 1}

All heads are nominative. The English gloss states the essential meaning; the accompanying note identifies the principal target or a likely error.

\par\noindent\begin{minipage}{\linewidth}\raggedright
% item: EX-I-01-01
% key: KEY-EX-I-01-01
\itempara{1}{\textbf{Ðæt ǣfen and se morgen.} — The evening and the morning. Neuter evening contrasts with masculine morning; both are singular.}

\end{minipage}\par

\par\noindent\begin{minipage}{\linewidth}\raggedright
% item: EX-I-01-02
% key: KEY-EX-I-01-02
\itempara{2}{\textbf{Se dæg and sēo niht.} — The day and the night. The article distinguishes masculine day and feminine night.}

\end{minipage}\par

\par\noindent\begin{minipage}{\linewidth}\raggedright
% item: EX-I-01-03
% key: KEY-EX-I-01-03
\itempara{3}{\textbf{Sēo eorðe and se heofon.} — The earth and the heaven. Use weak feminine eorðe and the selected strong masculine heofon.}

\end{minipage}\par

\par\noindent\begin{minipage}{\linewidth}\raggedright
% item: EX-I-01-04
% key: KEY-EX-I-01-04
\itempara{4}{\textbf{Se fisc and se fugol.} — The fish and the bird. Both heads are masculine nominative singular.}

\end{minipage}\par

\par\noindent\begin{minipage}{\linewidth}\raggedright
% item: EX-I-01-05
% key: KEY-EX-I-01-05
\itempara{5}{\textbf{Ðæt gærs and ðæt sǣd.} — The grass and the seed. Both are neuter. Retain the source spelling gærs.}

\end{minipage}\par

\par\noindent\begin{minipage}{\linewidth}\raggedright
% item: EX-I-01-06
% key: KEY-EX-I-01-06
\itempara{6}{\textbf{Ðæt trēo and sēo wyrt.} — The tree and the herb. The tree is neuter, the herb feminine; trēo is the bare singular.}

\end{minipage}\par

\par\noindent\begin{minipage}{\linewidth}\raggedright
% item: EX-I-01-07
% key: KEY-EX-I-01-07
\itempara{7}{\textbf{Sēo drȳgnes and ðæt wæter.} — The dry land and the water. Dry land is feminine drȳgnes; water is neuter. No ending is added to either nominative.}

\end{minipage}\par

\par\noindent\begin{minipage}{\linewidth}\raggedright
% item: EX-I-01-08
% key: KEY-EX-I-01-08
\itempara{8}{\textbf{Sēo andlīcnes and sēo gelīcnes.} — The image and the likeness. Both are feminine; the doubled s belongs to their inflected stems, not these nominatives.}

\end{minipage}\par

\par\noindent\begin{minipage}{\linewidth}\raggedright
% item: EX-I-01-09
% key: KEY-EX-I-01-09
\itempara{9}{\textbf{Ðæt lēoht and ðā ðīestru.} — The light and the darkness (plur.). The light is singular; darkness is explicitly plural, so its article is ðā.}

\end{minipage}\par

\par\noindent\begin{minipage}{\linewidth}\raggedright
% item: EX-I-01-10
% key: KEY-EX-I-01-10
\itempara{10}{\textbf{Ðæt tācen and sēo tīd.} — The sign and the season. The sign is neuter; the season is feminine. Tācen retains its e in the bare singular.}

\end{minipage}\par

\par\noindent\begin{minipage}{\linewidth}\raggedright
% item: EX-I-01-11
% key: KEY-EX-I-01-11
\itempara{11}{\textbf{Se gāst and ðæt līf.} — The spirit and the life. Masculine spirit contrasts with neuter life.}

\end{minipage}\par

\par\noindent\begin{minipage}{\linewidth}\raggedright
% item: EX-I-01-12
% key: KEY-EX-I-01-12
\itempara{12}{\textbf{Sēo gesceaft and se męte.} — The creature and the food. The selected creature is feminine, food masculine; preserve hooked ę.}

\end{minipage}\par

\par\noindent\begin{minipage}{\linewidth}\raggedright
% item: EX-I-01-13
% key: KEY-EX-I-01-13
\itempara{13}{\textbf{Ðæt gēar and se dæg.} — The year and the day. Year is neuter, day masculine. The dæg stem changes only in the plural.}

\end{minipage}\par

\par\noindent\begin{minipage}{\linewidth}\raggedright
% item: EX-I-01-14
% key: KEY-EX-I-01-14
\itempara{14}{\textbf{Se mann and ðæt dēor.} — The man and the animal. The selected animal is neuter dēor; singular mann has no vowel change.}

\end{minipage}\par

\par\noindent\begin{minipage}{\linewidth}\raggedright
% item: EX-I-01-15
% key: KEY-EX-I-01-15
\itempara{15}{\textbf{Ðæt trēo and se wæstm.} — The tree and the fruit. Fruit is masculine wæstm, with short æ. The tree is singular.}

\end{minipage}\par

\par\noindent\begin{minipage}{\linewidth}\raggedright
% item: EX-I-01-16
% key: KEY-EX-I-01-16
\itempara{16}{\textbf{Sēo fæstnes and sēo nēowolnes.} — The firmament and the abyss. Both are feminine nominatives in -nes; the assigned senses are firmament and abyss.}

\end{minipage}\par

\subsection*{Group I, Exercise 2}

All noun phrases are plural. In particular, zero-ending neuters and \textbf{niht} must not be made singular simply because their spelling is unchanged. The plural article \textbf{ðā} is common to every gender.

\par\noindent\begin{minipage}{\linewidth}\raggedright
% item: EX-I-02-01
% key: KEY-EX-I-02-01
\itempara{1}{\textbf{Ðā dagas and ðā niht.} — The days and the nights. Dæg has plural a; niht has an endingless N/A plural.}

\end{minipage}\par

\par\noindent\begin{minipage}{\linewidth}\raggedright
% item: EX-I-02-02
% key: KEY-EX-I-02-02
\itempara{2}{\textbf{Ðā fiscas and ðā fuglas.} — The fishes and the birds. Both have masculine plural -as; fugol uses the selected syncopated stem fugl-.}

\end{minipage}\par

\par\noindent\begin{minipage}{\linewidth}\raggedright
% item: EX-I-02-03
% key: KEY-EX-I-02-03
\itempara{3}{\textbf{Ðā trēowu and ðā wyrta.} — The trees and the herbs. Trēowu is the exceptional neuter plural; wyrta is feminine N/A plural.}

\end{minipage}\par

\par\noindent\begin{minipage}{\linewidth}\raggedright
% item: EX-I-02-04
% key: KEY-EX-I-02-04
\itempara{4}{\textbf{Ðā andlīcnessa and ðā gelīcnessa.} — The images and the likenesses. Both -nes nouns have -nessa. Single s before -a misses the selected stem.}

\end{minipage}\par

\par\noindent\begin{minipage}{\linewidth}\raggedright
% item: EX-I-02-05
% key: KEY-EX-I-02-05
\itempara{5}{\textbf{Ðā tācnu and ðā tīda.} — The signs and the seasons. Tācen loses e before the plural ending; tīd takes the selected feminine -a.}

\end{minipage}\par

\par\noindent\begin{minipage}{\linewidth}\raggedright
% item: EX-I-02-06
% key: KEY-EX-I-02-06
\itempara{6}{\textbf{Ðā gēar and ðā dagas.} — The years and the days. Gēar has zero N/A plural, while dæg changes its vowel and takes -as.}

\end{minipage}\par

\par\noindent\begin{minipage}{\linewidth}\raggedright
% item: EX-I-02-07
% key: KEY-EX-I-02-07
\itempara{7}{\textbf{Ðā męnn and ðā dēor.} — The men and the animals. Męnn is the changed-vowel plural; dēor has no N/A plural ending.}

\end{minipage}\par

\par\noindent\begin{minipage}{\linewidth}\raggedright
% item: EX-I-02-08
% key: KEY-EX-I-02-08
\itempara{8}{\textbf{Ðā sǣs and ðā heofonas.} — The seas and the heavens. Sǣs is the irregular sea plural; heofonas retains the middle vowel.}

\end{minipage}\par

\par\noindent\begin{minipage}{\linewidth}\raggedright
% item: EX-I-02-09
% key: KEY-EX-I-02-09
\itempara{9}{\textbf{Ðā sǣd and ðā wæstmas.} — The seeds and the fruits. Neuter seeds have zero ending; masculine fruits take -as without changing æ.}

\end{minipage}\par

\par\noindent\begin{minipage}{\linewidth}\raggedright
% item: EX-I-02-10
% key: KEY-EX-I-02-10
\itempara{10}{\textbf{Ðā ǣfenu and ðā morgnas.} — The evenings and the mornings. The selected neuter ǣfenu retains e; selected morgnas loses it. Morgenas is also valid.}

\end{minipage}\par

\par\noindent\begin{minipage}{\linewidth}\raggedright
% item: EX-I-02-11
% key: KEY-EX-I-02-11
\itempara{11}{\textbf{Ðā gāstas and ðā męnn.} — The spirits and the men. Gāst has regular -as; mann requires męnn, not an -as plural.}

\end{minipage}\par

\par\noindent\begin{minipage}{\linewidth}\raggedright
% item: EX-I-02-12
% key: KEY-EX-I-02-12
\itempara{12}{\textbf{Ðā hwalas and ðā nīetenu.} — The whales and the cattle. Hwæl changes æ to a; nīetenu retains the middle e and supplies the assigned cattle sense.}

\end{minipage}\par

\subsection*{Group I, Exercise 3}

The selected target is \textbf{genitive dependent + nominative head}, with no separate article on the head. The article of the dependent remains. This is Cook's preferred pattern for this exercise, not a ban on every other grammatical genitive order. Proper \textbf{Godes} has no article. A plural head and a plural dependent are independent choices: one does not force the other to be plural.

\par\noindent\begin{minipage}{\linewidth}\raggedright
% item: EX-I-03-01
% key: KEY-EX-I-03-01
\itempara{1}{\textbf{Ðæs dæges anginn.} — The beginning of the day. Day is genitive singular; beginning is the neuter nominative head.}

\end{minipage}\par

\par\noindent\begin{minipage}{\linewidth}\raggedright
% item: EX-I-03-02
% key: KEY-EX-I-03-02
\itempara{2}{\textbf{Godes gāst.} — The spirit of God. The proper divine name has genitive Godes without an article; gāst is nominative.}

\end{minipage}\par

\par\noindent\begin{minipage}{\linewidth}\raggedright
% item: EX-I-03-03
% key: KEY-EX-I-03-03
\itempara{3}{\textbf{Ðæs līfes lēoht.} — The light of the life. Life is genitive singular. Its article is retained, while the head light has none of its own.}

\end{minipage}\par

\par\noindent\begin{minipage}{\linewidth}\raggedright
% item: EX-I-03-04
% key: KEY-EX-I-03-04
\itempara{4}{\textbf{Ðæs fugles līf.} — The life of the bird. Fugles is the selected singular genitive with syncope; līf is the nominative head.}

\end{minipage}\par

\par\noindent\begin{minipage}{\linewidth}\raggedright
% item: EX-I-03-05
% key: KEY-EX-I-03-05
\itempara{5}{\textbf{Ðæs fisces męte.} — The food of the fish. Fish is genitive singular; food is nominative męte, not a second genitive.}

\end{minipage}\par

\par\noindent\begin{minipage}{\linewidth}\raggedright
% item: EX-I-03-06
% key: KEY-EX-I-03-06
\itempara{6}{\textbf{Ðæs heofones fæstnes.} — The firmament of the heaven. Heofones retains its e. The head fæstnes is nominative, without -e.}

\end{minipage}\par

\par\noindent\begin{minipage}{\linewidth}\raggedright
% item: EX-I-03-07
% key: KEY-EX-I-03-07
\itempara{7}{\textbf{Ðǣre eorðan dēor.} — The beasts of the earth. Earth is singular genitive; beasts are plural dēor, though the head has no plural ending.}

\end{minipage}\par

\par\noindent\begin{minipage}{\linewidth}\raggedright
% item: EX-I-03-08
% key: KEY-EX-I-03-08
\itempara{8}{\textbf{Ðǣre nēowolnesse brādnes.} — The face of the deep. The deep is genitive singular; brādnes is the nominative head, meaning surface/face here, not a human face.}

\end{minipage}\par

\par\noindent\begin{minipage}{\linewidth}\raggedright
% item: EX-I-03-09
% key: KEY-EX-I-03-09
\itempara{9}{\textbf{Ðāra wætera gegaderunga.} — The congregations of the waters. Both waters and gatherings are plural: genitive wætera and nominative gegaderunga.}

\end{minipage}\par

\par\noindent\begin{minipage}{\linewidth}\raggedright
% item: EX-I-03-10
% key: KEY-EX-I-03-10
\itempara{10}{\textbf{Ðæs sǣs wæteru.} — The waters of the sea. The sea is singular genitive sǣs; waters are nominative plural wæteru. No sea genitive plural is needed.}

\end{minipage}\par

\par\noindent\begin{minipage}{\linewidth}\raggedright
% item: EX-I-03-11
% key: KEY-EX-I-03-11
\itempara{11}{\textbf{Ðæs sǣs fiscas.} — The fishes of the sea. Sea remains genitive singular, while fishes is the masculine nominative plural.}

\end{minipage}\par

\par\noindent\begin{minipage}{\linewidth}\raggedright
% item: EX-I-03-12
% key: KEY-EX-I-03-12
\itempara{12}{\textbf{Ðæs lyftes fuglas.} — The birds of the air. Use the selected neuter air, hence ðæs lyftes; birds are plural fuglas.}

\end{minipage}\par

\par\noindent\begin{minipage}{\linewidth}\raggedright
% item: EX-I-03-13
% key: KEY-EX-I-03-13
\itempara{13}{\textbf{Ðāra heofona and ðǣre eorðan frætwung.} — The array of the heavens and of the earth. Two coordinated genitives modify one singular head: heavens plural, earth singular. Each dependent keeps its own article.}

\end{minipage}\par

\par\noindent\begin{minipage}{\linewidth}\raggedright
% item: EX-I-03-14
% key: KEY-EX-I-03-14
\itempara{14}{\textbf{Godes weorc.} — The works of God. Godes is singular; weorc is nominative plural, with no added ending.}

\end{minipage}\par

\par\noindent\begin{minipage}{\linewidth}\raggedright
% item: EX-I-03-15
% key: KEY-EX-I-03-15
\itempara{15}{\textbf{Ðāra dēora cynn.} — The kinds of the animals. The animals are genitive plural; the kinds are nominative plural cynn, again with no added ending.}

\end{minipage}\par

\par\noindent\begin{minipage}{\linewidth}\raggedright
% item: EX-I-03-16
% key: KEY-EX-I-03-16
\itempara{16}{\textbf{Ðæs heofones steorran.} — The stars of the heaven. Heaven is singular genitive; stars are weak masculine nominative plural -an.}

\end{minipage}\par

\par\noindent\begin{minipage}{\linewidth}\raggedright
% item: EX-I-03-17
% key: KEY-EX-I-03-17
\itempara{17}{\textbf{Ðæs gærses sǣd.} — The seeds of the grass. Grass is singular genitive. Seeds are plural sǣd, despite the unchanged spelling.}

\end{minipage}\par

\par\noindent\begin{minipage}{\linewidth}\raggedright
% item: EX-I-03-18
% key: KEY-EX-I-03-18
\itempara{18}{\textbf{Ðāra nihta līehting.} — The illumination of the nights. Nights are genitive plural nihta; illumination is the feminine nominative singular head. Retain the macron of līehting.}

\end{minipage}\par

\section*{14. Composition}

\par\noindent\begin{minipage}{\linewidth}\raggedright
% item: CH01-COMP01
% key: KEY-CH01-COMP01
\itempara{1}{\textbf{Se fugol is on ðǣm trēowe.} — The bird is in the tree. Singular subject takes is; trēowe is dative after locative on, with the required w.}

\end{minipage}\par

\par\noindent\begin{minipage}{\linewidth}\raggedright
% item: CH01-COMP02
% key: KEY-CH01-COMP02
\itempara{2}{\textbf{Ðā fuglas wǣron on ðǣm trēowum.} — The birds were in the trees. Plural past wǣron agrees with fuglas; trēowum is dative plural.}

\end{minipage}\par

\par\noindent\begin{minipage}{\linewidth}\raggedright
% item: CH01-COMP03
% key: KEY-CH01-COMP03
\itempara{3}{\textbf{Ðæt sǣd is ðæs fugles męte.} — The seed is the bird’s food. Męte is a nominative noun predicate; ðæs fugles is its dependent genitive. Do not add a separate article to męte in the requested pattern.}

\end{minipage}\par

\par\noindent\begin{minipage}{\linewidth}\raggedright
% item: CH01-COMP04
% key: KEY-CH01-COMP04
\itempara{4}{\textbf{Se fisc is dēor.} — The fish is an animal. Dēor is a nominative predicate; English an needs no separate Old English word here.}

\end{minipage}\par

\par\noindent\begin{minipage}{\linewidth}\raggedright
% item: CH01-COMP05
% key: KEY-CH01-COMP05
\itempara{5}{\textbf{Ðā tācnu wǣron on ðǣm heofone.} — The signs were in the sky. Tācnu is nominative plural with syncope; heofone is dative singular with retained e.}

\end{minipage}\par

\par\noindent\begin{minipage}{\linewidth}\raggedright
% item: CH01-COMP06
% key: KEY-CH01-COMP06
\itempara{6}{\textbf{Ðā wyrta sindon on ðæs mannes hofe.} — The herbs are in the man’s dwelling. Plural subject takes sindon. The genitive dependent does not replace the head’s dative after on.}

\end{minipage}\par

\par\noindent\begin{minipage}{\linewidth}\raggedright
% item: CH01-COMP07
% key: KEY-CH01-COMP07
\itempara{7}{\textbf{Ðǣre niht lēoht wæs on ðǣm heofone.} — The light of the night was in the sky. Niht is selected genitive singular, despite its lack of ending; ðǣre agrees with genitive niht here. The singular head lēoht takes wæs. Source-listed nihte is also a valid genitive alternative.}

\end{minipage}\par

\par\noindent\begin{minipage}{\linewidth}\raggedright
% item: CH01-COMP08
% key: KEY-CH01-COMP08
\itempara{8}{\textbf{Ðāra dēora męte is ðāra trēowa wæstm.} — The animals’ food is the fruit of the trees. Both dependents are genitive plural. Subject męte and predicate wæstm are singular nominatives, so the copula is is. Trēowa retains w, and wæstm retains short æ.}

\end{minipage}\par

\section*{15. Short readings}

\par\noindent\begin{minipage}{\linewidth}\raggedright
% item: CH01-MINI01
% key: KEY-CH01-MINI01
\itempara{1}{\textbf{The man is in the dwelling. The seed is in the man's dwelling. The bird is in the tree. The seed is the bird's food.} The dwelling in the second sentence belongs to the man. \textbf{Mannes} is genitive singular, with its agreeing \textbf{ðæs}; \textbf{hofe} is dative singular, governed by locative \textbf{on}. In the last sentence, \textbf{męte} identifies the subject and is nominative, while \textbf{ðæs fugles} is its genitive dependent. “The food of the bird” conveys the same relation.}

\end{minipage}\par

\par\noindent\begin{minipage}{\linewidth}\raggedright
% item: CH01-MINI02
% key: KEY-CH01-MINI02
\itempara{2}{\textbf{The fishes were in the sea. The birds were in the trees. The trees were on the earth.} “Fish” as the English plural is equally suitable.}

\continuation{Present: \textbf{Ðā fiscas sindon on ðǣm sǣ. Ðā fuglas sindon on ðǣm trēowum. Ðā trēowu sindon on ðǣre eorðan.}}

\continuation{\textbf{Trēowu} is nominative plural, the subject of the third sentence; \textbf{trēowum} in the second is dative plural after \textbf{on}. \textbf{Eorðan} is dative singular after locative \textbf{on}, as its feminine singular article \textbf{ðǣre} confirms. Replacing \textbf{wǣron} with singular \textbf{is} would lose agreement.}

\end{minipage}\par

\section*{16. The first day of Creation}

\par\noindent\begin{minipage}{\linewidth}\raggedright
% item: CH01-READ01
% key: KEY-CH01-READ01
\itempara{1}{\textbf{God} is the subject; \textbf{heofonan} and \textbf{eorðan} are the accusative objects. \textbf{Heofonan} is singular, the weak feminine form supplied in the help. \textbf{In the beginning God created heaven and earth.} The opening places the verb before its subject after \textbf{On anginne}. Do not identify \textbf{God} as an object simply from its position, or translate \textbf{heofonan} as a plural of the masculine \textbf{heofon}. English “the heaven(s) and the earth” may be idiomatic, but the grammatical analysis must retain the source's singular.}

\end{minipage}\par

\par\noindent\begin{minipage}{\linewidth}\raggedright
% item: CH01-READ02
% key: KEY-CH01-READ02
\itempara{2}{\textbf{The earth was empty and void; darkness was over the surface of the deep.} \textbf{Ðīestru} is plural and therefore has plural \textbf{wǣron}, although ordinary English uses singular “darkness was.” After \textbf{ofer}, the head is \textbf{brādnesse}, dative singular. Its dependent \textbf{ðǣre nēowolnesse} is genitive singular. The article belongs to the dependent. Treating every word after \textbf{ofer} as dative misses the internal structure of the phrase. “Abyss” for “deep” and “face” for “surface” are defensible.}

\end{minipage}\par

\par\noindent\begin{minipage}{\linewidth}\raggedright
% item: CH01-READ03
% key: KEY-CH01-READ03
\itempara{3}{\textbf{God's spirit was borne/carried over the waters}, or contextually \textbf{moved over the waters}. \textbf{Wæs gefęred} is the supplied passive expression. \textbf{Godes} is the genitive dependent of nominative \textbf{gāst}. \textbf{Wæteru} is accusative plural here; the selected dative plural would be \textbf{wæterum}. The previous \textbf{ofer + dative} does not make this phrase dative too. A translation asserting that the spirit carried some unexpressed object changes the construction.}

\end{minipage}\par

\par\noindent\begin{minipage}{\linewidth}\raggedright
% item: CH01-READ04
% key: KEY-CH01-READ04
\itempara{4}{The speech is a \textbf{command or wish for light to come into being}, not a report that light already exists: \textbf{“Let light come to be.”} It is followed by \textbf{light was made / came into being}. “Let there be light” gives the intended force naturally. The subsequent \textbf{wearð geworht} is the supplied passive; do not change the quotation into indicative “light becomes.”}

\end{minipage}\par

\par\noindent\begin{minipage}{\linewidth}\raggedright
% item: CH01-READ05
% key: KEY-CH01-READ05
\itempara{5}{\textbf{Ðā} is temporal “then”; \textbf{ðæt} introduces the content clause, “that”; \textbf{hit} is nominative neuter singular “it,” referring to \textbf{lēoht}. \textbf{God then saw that it was good.} The final \textbf{wæs} belongs to the embedded clause, whose subject is \textbf{hit}. This \textbf{ðæt} is not the article; \textbf{gōd}, “good,” is not the proper name \textbf{God}.}

\end{minipage}\par

\par\noindent\begin{minipage}{\linewidth}\raggedright
% item: CH01-READ06
% key: KEY-CH01-READ06
\itempara{6}{God separated \textbf{the light from the darkness}. \textbf{Ðæt lēoht} is the accusative singular direct object; \textbf{fram ðǣm ðīestrum} is the separation phrase, with dative plural after \textbf{fram}. \textbf{Hē} refers to God. “Divided” is a suitable alternative to “separated.” Reversing the object and the phrase loses the stated construction.}

\end{minipage}\par

\par\noindent\begin{minipage}{\linewidth}\raggedright
% item: CH01-READ07
% key: KEY-CH01-READ07
\itempara{7}{He named \textbf{the light Day} and \textbf{the darkness Night}. \textbf{God/he} is the understood subject of \textbf{hēt}. The ending states \textbf{one day}. The naming objects and the names applied to them are accusative. A translation may supply the understood English “he.” \textbf{Ān} is cardinal “one,” so “the first day” is not an exact rendering of the closing words. The supplied closing formula is sufficient; no unique analysis of its subject/predicate arrangement is required.}

\end{minipage}\par

\par\noindent\begin{minipage}{\linewidth}\raggedright
% item: CH01-READ08
% key: KEY-CH01-READ08
\itempara{8}{\textbf{In the beginning God created heaven and earth. Now the earth was empty and void; and darkness was over the surface of the deep; and God's spirit was borne over the waters. God then said, “Let light come to be”; and light was made. God then saw that it was good; and he separated the light from the darkness. And he called the light Day, and the darkness Night. Then there came to be evening and morning: one day.}}

\continuation{The translation follows the actual Old English. “Moved” is natural for \textbf{wæs gefęred} in context; “came into being” is suitable for \textbf{wearð geworht}. In the close, “evening and morning had come to be: one day” preserves the resultative/perfect force of \textbf{wæs geworden}. The latter is a coming-to-be expression, not simply another instance of the earlier passive. The connective \textbf{sōðlice} may be rendered “now,” “for its part,” or by a suitable discourse connection. Preserve creation, description, movement, command, fulfillment, judgment, separation, naming, and the one-day close. A remembered translation of Genesis is not sufficient if it substitutes different grammar or wording for this extract.}

\end{minipage}\par

\section*{17. Review}

\par\noindent\begin{minipage}{\linewidth}\raggedright
% item: CH01-REVIEW01
% key: KEY-CH01-REVIEW01
\itempara{1}{\textbf{Dagas, hwalas, wæstmas.} The first two nouns have \textbf{æ} in the singular and \textbf{a} in the plural. \textbf{Wæstm} retains short \textbf{æ} before its consonant group; its plural adds \textbf{-as} without that vowel change. The correction of Cook's length mark does not authorize changing the vowel to \textbf{a}. The target is controlled use of the particular stem alternation, not a universal replacement of \textbf{æ}.}

\end{minipage}\par

\par\noindent\begin{minipage}{\linewidth}\raggedright
% item: CH01-REVIEW02
% key: KEY-CH01-REVIEW02
\itempara{2}{\textbf{Hofe} is the head, dative singular after locative \textbf{on}. \textbf{Ðæs mannes} is the genitive singular dependent, “the man's.” \textbf{Ðæs} agrees with \textbf{mannes}. The head's externally governed dative and the dependent's internal genitive coexist; they need not have the same case.}

\end{minipage}\par

\par\noindent\begin{minipage}{\linewidth}\raggedright
% item: CH01-REVIEW03
% key: KEY-CH01-REVIEW03
\itempara{3}{\textbf{Eorðan} can be genitive, dative, or accusative singular, or nominative/accusative plural, in the displayed weak paradigm. In \textbf{on ðǣre eorðan}, it is \textbf{dative singular}, with the locative preposition and feminine singular article. In \textbf{God gescēop ðā eorðan}, it is \textbf{accusative singular}: the direct-object construction establishes the case, and the supplied meaning “God created the earth” establishes the singular number. The spellings \textbf{ðā/eorðan} alone also permit \textbf{accusative plural}; they do not by themselves establish singular number. Naming the possible plural cells does not assign a new plural sense of “earth” for composition.}

\end{minipage}\par

\par\noindent\begin{minipage}{\linewidth}\raggedright
% item: CH01-REVIEW04
% key: KEY-CH01-REVIEW04
\itempara{4}{\textbf{Ðīestru} is a plural noun, so the Old English verb is plural \textbf{wǣron}; English “darkness” is ordinarily singular. In \textbf{ðæt lēoht fram ðǣm ðīestrum}, \textbf{ðæt lēoht} is accusative singular, the thing separated; \textbf{ðǣm ðīestrum} is dative plural after \textbf{fram}. Translating English number directly into the Old English verb would give the wrong agreement.}

\end{minipage}\par

\clearpage
\section*{Source notes}

\begingroup\fontsize{10}{12.8}\selectfont

These are new models, not a key supplied by Cook. Source prompts are Cook, \emph{Exercises in Old English} (1895), Group I, Exercises 1–3, pp.1–2 / PDF9–10. The complete reading is \emph{A First Book in Old English}, second edition (1895), Reader I, pp.123–124 / PDF139–140. The learner's source notes give the checked grammatical and lexical locators for the models and their alternatives. PDF references are one-based. The reading translation was prepared from the printed Old English; no inherited teaching or previous-course answer serves as a prerequisite.

\endgroup
\end{document}
